\documentclass[12pt]{article}
\usepackage{amsmath,amssymb,amsthm}
\usepackage[margin=1in]{geometry}

\newtheorem{theorem}{Theorem}
\newtheorem{lemma}[theorem]{Lemma}

\title{Problem 492: Totient Chains}
\author{}
\date{}

\begin{document}
\maketitle

\section{Problem Statement}
Define the iterated Euler totient by $\varphi^{(0)}(n) = n$ and $\varphi^{(k)}(n) = \varphi(\varphi^{(k-1)}(n))$. The \emph{totient chain} starting at $n$ is $n \to \varphi(n) \to \varphi^{(2)}(n) \to \cdots \to 1$. Let $L(n)$ denote its length. Compute $\sum_{n=2}^{N} L(n)$.

\section{Mathematical Foundation}

\begin{theorem}[Strict decrease]
For all $n \ge 2$, $\varphi(n) < n$.
\end{theorem}

\begin{proof}
Since $n \ge 2$ has at least one prime factor $p$, the Euler product formula gives $\varphi(n) = n\prod_{p \mid n}(1 - 1/p) < n$ because each factor $(1 - 1/p) < 1$.
\end{proof}

\begin{theorem}[Well-definedness of $L$]
For every $n \ge 1$, the totient chain reaches $1$ in finitely many steps. Moreover, $L(1) = 0$ and $L(n) = 1 + L(\varphi(n))$ for $n \ge 2$.
\end{theorem}

\begin{proof}
By Theorem~1, the chain is a strictly decreasing sequence of positive integers for all terms $\ge 2$. Being bounded below by $1$, it terminates. The recurrence is immediate from the definition.
\end{proof}

\begin{lemma}[Logarithmic bound]
$L(n) = O(\log n)$ for all $n \ge 1$.
\end{lemma}

\begin{proof}
For $n \ge 3$, $\varphi(n)$ is even. For even $m \ge 2$, $\varphi(m) \le m/2$. After at most one step, each subsequent application at least halves the argument, so $L(n) \le 1 + \log_2 n + 1 = O(\log n)$.
\end{proof}

\begin{theorem}[Correctness of Euler product sieve]
Initializing $\varphi[n] = n$ for all $n$ and, for each prime $p$, setting $\varphi[kp] \leftarrow \varphi[kp] \cdot (p-1)/p$ for all multiples $kp$, correctly computes $\varphi(n) = n\prod_{p \mid n}(1 - 1/p)$.
\end{theorem}

\begin{proof}
Each prime $p$ is processed exactly once. For every multiple $n$ of $p$, the factor $(1 - 1/p)$ is multiplied in. After all primes up to $N$ are processed, $\varphi[n]$ equals the Euler product.
\end{proof}

\section{Algorithm}
\begin{enumerate}
\item Compute $\varphi(n)$ for $n = 1, \ldots, N$ via the Euler product sieve.
\item Compute $L(n)$ bottom-up: $L(1) = 0$; for $n = 2, \ldots, N$, set $L(n) = 1 + L(\varphi(n))$.
\item Accumulate $\sum_{n=2}^{N} L(n)$.
\end{enumerate}

\section{Complexity Analysis}
\begin{itemize}
\item \textbf{Time:} $O(N \log\log N)$ for the sieve $+$ $O(N)$ for chain lengths $= O(N \log\log N)$.
\item \textbf{Space:} $O(N)$ for the arrays $\varphi[\cdot]$ and $L[\cdot]$.
\end{itemize}

\section{Answer}

\[
\boxed{242586962923928}
\]

\end{document}
